\documentclass[a4paper,12pt]{article}
\usepackage[utf8]{inputenc}
\usepackage[T1]{fontenc}
\usepackage[margin=2.5cm]{geometry}
\usepackage{ctex} % 支持中文
\usepackage{setspace} % 设置行距
\usepackage{titlesec} % 标题格式
\usepackage{xcolor} % 颜色支持
\usepackage{enumitem} % 列表优化

% 设置行距为1.5倍，符合标准商务报告阅读习惯
\onehalfspacing

% 设置标题格式
\titleformat{\section}{\Large\bfseries\color{blue!40!black}}{}{0.5em}{}
\titleformat{\subsection}{\large\bfseries\color{gray!80!black}}{}{0.5em}{}

% 列表设置
\setlist[itemize]{leftmargin=*}

% 元数据
\title{\textbf{2024年度环境、社会及管治（ESG）报告节选}}
\author{公司董事会办公室}
\date{2024年3月}

\begin{document}

\maketitle

% 根据指令，不包含摘要和目录

\section{1. 前言：战略定位与可持续发展承诺}

2024年是海南自贸港封关运作准备的关键之年。作为深度参与自贸港建设的大型基础设施及临空产业上市公司，本公司（以下简称“公司”）紧密围绕“一主一次，两支一货运”的区域机场群布局，积极履行企业社会责任。公司致力于通过高标准的治理架构、环境管理体系及社会价值创造，服务国家开放战略，并为股东及利益相关方创造长期价值。

\section{2. 环境范畴（Environmental）：低碳运营与资源管理}

公司严格遵守国家及地方法律法规，制定并执行环境管理策略。2024年，公司持续推进“双碳”目标落地，在能源管理、温室气体减排及绿色建筑认证方面取得实质性进展，确保环境绩效的可测量与可追溯。

\subsection{2.1 排放管理与资源消耗绩效}

报告期内，公司完成了年度温室气体排放核查及能源审计工作。根据核算结果，2024年度公司温室气体排放总量为 \textbf{51,518.71} 吨二氧化碳当量（tCO$_2$e）。其中，源自化石燃料燃烧的直接排放（范围1）为 \textbf{11,214.39} tCO$_2$e，源自外购电力的间接排放（范围2）为 \textbf{40,304.32} tCO$_2$e。经计算，公司年度碳排放强度为 0.12 吨二氧化碳当量/万元，保持在行业合理区间。

在能源与资源消耗方面，全年综合能源消耗量折合标准煤 \textbf{13,538.22} 吨。具体能源种类消耗显示，外购电力消耗量为 7,511.05 万千瓦时，汽油消耗量为 19.10 万升。水资源管理方面，公司全年水资源消耗总量为 145.94 万吨，通过不断优化的中水回收系统，实现中水回用量 13.61 万吨，有效降低了新水取用量。此外，公司严格管控废弃物排放，全年产生无害废弃物 6,230.58 吨，有害废弃物 0.58 吨，所有废弃物均已委托有资质单位进行处置，合规处理率达到 100\%。

\subsection{2.2 绿色机场建设与减排措施}

为响应民航局“绿色机场”建设号召，公司重点推进了辅助动力装置（APU）替代设施的应用与清洁能源建设。

在APU替代设施建设方面，公司在旗下核心门户机场全面推广使用地面电源和空调设备替代飞机辅助动力装置，以减少航空器在地面运行期间的燃油消耗及噪音污染。截至报告期末，已投入使用APU替代设备共计 \textbf{31套}（含近机位18套，远机位13套）。经测算，该项措施全年累计减少碳排放 \textbf{25,382} 吨二氧化碳当量。凭借在减排领域的突出表现，旗下核心门户机场成功获评“双碳机场”二星级称号。

在可再生能源利用方面，公司积极利用自有资产空间布局光伏发电项目。旗下核心商业综合体屋顶分布式光伏项目运行稳定，年均发电量达 450 万千瓦时，约占该商业体总用电量的 10\%，折合年减排量 2,263.59 吨二氧化碳。此外，旗下临空产业园利用厂房屋顶完成光伏装机容量 8.6MWp，预计年减排贡献可达 25.75 万吨。与此同时，公司持续推进场内车辆“油改电”进程，核心机场新能源车辆保有量已达 221 辆，占比提升至 43.6\%，并配套新建充电桩约 100 个，车辆电动化全年实现减碳 761 吨。

\subsection{2.3 绿色建筑认证}

公司在项目开发建设全生命周期中融入绿色理念，严格遵循绿色建筑标准。省会城市在建地标项目“海南中心”（双子塔）已获得 \textbf{LEED CS 金级预认证}，体现了公司在超高层建筑节能设计上的国际化水准；某新区F07项目亦通过绿色建筑一星级标准认证。

\section{3. 社会范畴（Social）：运营韧性与共享价值}

公司坚持“真情服务”底线，持续提升安全保障能力与服务品质，注重员工权益保护，积极投身社会公益事业，致力于构建和谐的空港社区。

\subsection{3.1 运营绩效与市场表现}

2024年，得益于旅游市场的强劲复苏及航线网络的持续优化，公司核心业务指标实现稳步增长。报告期内，旗下控股及管理输出的9家机场总体完成起降 \textbf{16.11 万架次}，旅客吞吐量共计 \textbf{2,462.62 万人次}。

其中，旗下核心门户机场旅客吞吐量重回“2000万级”梯队，位列全国第29位，枢纽能级进一步巩固。旗下重点区域机场积极拓展国际网络，成功开通至吉隆坡的首条国际客运航线，提升了区域对外开放水平。

\subsection{3.2 安全管理与应急响应体系}

公司坚持“安全第一”原则，持续完善安全管理体系（SMS），重点强化极端天气应对能力及数字化安全管控。

在应对重大气象灾害方面，面对2024年9月超强台风“摩羯”的正面冲击，公司展现了高效的应急响应能力。台风登陆前，公司迅速启动“防汛防风专项应急预案”，运管委统筹实施航班动态调整及航空器系留工作，确保资产安全。台风过境后，公司立即组织力量开展适航恢复，9月3日至7日期间，核心机场迅速恢复运行并累计保障航班 \textbf{1,300 架次}，有力保障了抢险救灾人员及物资的进出岛通道畅通，体现了基础设施运营者的责任担当。

在日常安全管理数字化转型方面，公司上线运行“安全智控”平台，实现了安全管理流程的标准化与数据化。该平台关联制度手册 9,626 项，全年流转电子检查单 31 万份，有效提升了风险隐患排查的及时性与准确性。报告期内，公司安全绩效表现优异，一般征候万架次率为 0，责任事故发生率为 0，安全培训覆盖率达到 100\%。

\subsection{3.3 客户服务与创新}

为提升旅客出行体验，公司重点推进了通关效率与支付便利化服务创新：
\begin{itemize}
    \item \textbf{通关流程优化}：实施“一机双屏”海关安检融合查验模式，实现旅客行李“一次过检”，显著缩短了旅客在核心控制区的排队等候时间。
    \item \textbf{支付便利化}：针对外籍来华人员支付痛点，在核心机场启用“一站式综合服务中心”，整合外卡支付绑定、通讯服务及文旅咨询功能，切实提升入境旅客的便利度。
\end{itemize}

\subsection{3.4 人力资源与社会贡献}

公司秉持多元包容的雇主品牌理念。截至2024年末，公司员工总数为 10,611 人。在员工多元化方面，女性员工占比为 38.9\%（4,133人），少数民族员工 179 人；董事会中女性成员占比 22\%。公司重视人才培养，全年培训总投入 1,119.54 万元，人均培训时长达到 105.88 小时。

在社会公益领域，公司积极响应乡村振兴战略，全年投入乡村振兴资金 \textbf{190 万元}。同时，公司鼓励员工参与志愿服务，全年志愿服务参与超 10,000 人次，累计时长达 500,000 小时。此外，公司圆满完成了博鳌亚洲论坛年会、消博会、少数民族运动会等国家级重大活动的保障任务。

\section{4. 管治范畴（Governance）：合规与风险控制}

公司致力于维持高水平的企业管治标准，通过完善的治理架构与监督体系，保障公司稳健运营及股东权益。

\subsection{4.1 治理架构与合规管理}

公司严格按照上市公司监管要求，构建了规范的“五会一层”治理体系。董事会由9名董事组成，其中独立董事3名，结构合理。报告期内，公司共召开股东大会5次，董事会13次，并发布公告70份，确保了信息披露的真实、准确、完整。

在内部控制方面，公司构建了“三全三化”大监督体系，全年梳理风险点 358 个并制定管控措施。公司高度重视商业道德建设，反腐败培训覆盖率在管理层达到 95\%，员工层达到 90\%，培训总时长 400 小时。供应链管理方面，公司要求所有供应商签署廉洁承诺书（签署率100\%），并对23家工程类供应商进行了严格遴选。

\subsection{4.2 经济绩效延伸}

在主营业务之外，公司依托自贸港政策优势，积极参与离岛免税商业布局。2024年，公司间接参与的5家离岛免税店实现线下销售额约 \textbf{52 亿元}，市场份额占比约 \textbf{17\%}，有效促进了区域消费经济的发展与回流。

\end{document}